\section{Conclusions and final summary}
\label{sec:conclusions}

Every conclusion below rests on the measured run (Section~\ref{sec:method}) and on
the code it was built from. Each one was checked against the data. Claims that did not survive that
check are not included.

\subsection{Conclusions supported by the data}

\paragraph{C1. A second, simplex connection per peer is not worth it.}
\begin{itemize}
  \item \textbf{No throughput gain.} At 32\,B it is slower than duplex: \mode{poll}
        158k against 188k msg/s, \mode{epoll} 169k against 189k. At 4 and 64\,KiB the
        two shapes are within noise.
  \item \textbf{A longer handshake.} It costs 3\,RTT instead of 2: 62 against 41\,ms at
        20\,ms RTT, and a p50 of about 280 against 95\,ms at 1,024 peers.
  \item \textbf{An unreliable handshake under churn.} It collapses from 40\,ms RTT
        under 16 reconnects/s, and it slows behind a rate limit (78 against 9\,ms).
  \item \textbf{Twice the file descriptors.}
\end{itemize}
The collapse comes from the dial-back being a blocking \code{connect()} on the shared
sender thread, not from the retry sleep. How abruptly it appears is partly set by the
benchmark's open-loop churn and its 90\%-ready check (Section~\ref{sec:res-wan}). A non-blocking dial-back would remove the cliff, but not the extra
round trip or the extra connection.

\paragraph{C2. Separate read and write pollers pay off when the single I/O thread is
the bottleneck.} The trigger can be a high message rate, large messages or heavy
fan-in; any one is enough.
\begin{itemize}
  \item \textbf{Large messages, few peers:} with only 8 peers at 4\,KiB, the split
        modes deliver 46\% and 76\% more than duplex.
  \item \textbf{Fan-in:} at 256 peers, p99 is 8.0--8.4\,ms against about 10.6\,ms, at
        lower CPU. At 1,024 peers, p50 is 60 against 76\,ms (\mode{poll}) and 147
        against 159\,ms (\mode{epoll}).
  \item \textbf{No gain} when the link is the bottleneck (S6b) or the load is mostly
        idle (S2).
  \item \textbf{Worse under churn:} more failed sends.
\end{itemize}

\paragraph{C3. \mode{peer-threads} is reliable and fast, but needs a thread cap, and
its retry path is a limitation of the current implementation.}
\begin{itemize}
  \item \textbf{Reliable:} all 96 of its cells are valid. It stays healthy in S6a at
        40--50\,ms RTT, where simplex collapses.
  \item \textbf{Fast:} 203k msg/s at 32\,B and 16.2k at 64\,KiB. It has the best p99
        at 256 peers (6.8\,ms) and shields other peers from a slow reader.
  \item \textbf{Needs a cap:} at 1,024 peers on two CPUs it delivers only 20.8k of
        102k msg/s. Capped at 256, it serves those peers at a p99 of 6.8\,ms and
        58\% CPU.
  \item \textbf{Limitation:} its writes are already non-blocking, but retries are a
        100\,ms sleep on the shared \code{SenderWorker} thread, which stalls every
        peer. That cost 29--33 failures, and a p99 of 64--103\,ms, at churn 50.
        \code{open()} also does a blocking \code{connect()} on the same thread.
\end{itemize}
The retry and the dial belong on the peer's own owner thread: a deadline in its poll
loop rather than a sleep, and an asynchronous dial. The shared sender would then
become a dispatcher into per-peer queues. Two caveats remain: \mode{peer-threads}
uses two fds per peer, and its backpressure path was never exercised (S3).

\paragraph{Also established.}
\begin{itemize}
  \item Thread-per-connection collapses at 1,024 peers on two CPUs, and below that
        ceiling it is the cheapest mode.
  \item \code{poll} pays for every idle socket; \code{epoll} does not.
  \item Only a blocking sender lets one slow reader stall everyone.
  \item Under churn, tail latency is set by the shared sender's retry sleep, not by
        the concurrency model or by isolating descriptors.
\end{itemize}

\paragraph{Which mode to use.}
\begin{itemize}
  \item \textbf{Up to a few hundred peers on few cores:} a thread per peer with
        pooled writes (\mode{peer-threads}, \mode{pool-sender}), with
        \code{-{}-max-peers} sized to the cores.
  \item \textbf{Many or mostly idle peers:} \mode{epoll-duplex}, or
        \mode{epoll-split} for large messages or high rates.
  \item \textbf{Churn over a WAN:} always duplex.
\end{itemize}

\subsection{Limitations and open observations}

\begin{itemize}
  \item The measurements come from one WSL2 host, count CPU for the node only, and
        give the node two CPUs, which sets where the thread ceiling falls.
  \item Two observations remain unexplained:
        \begin{itemize}
          \item why the poll-based duplex and split modes beat \code{epoll} at 1,024
                peers;
          \item why the duplex shapes never queued writes in S3.
        \end{itemize}
  \item The peer cap is not uniform across modes (S7).
  \item Open-loop churn turns sender overload into an abrupt exclusion rather than a
        gradual slowdown (S6a).
\end{itemize}

\subsection{Final summary: a hypothesis for a combined design}

The tested results above bracket a design we did not build or measure. We state it as
a hypothesis.

\paragraph{Proposal.} Keep the ownership model of \mode{peer-threads}, in which one
thread owns its sockets, reads them and reports their end. But give each owner thread
a \emph{small pool of peers} instead of one, and add a balancing policy. This is an
M:N reactor~\cite{reactor} that sits between the two extremes we measured: one peer
per pool (\mode{peer-threads}) and every peer in one pool (\mode{*-duplex} and
\mode{*-split}). Candidate balancing policies:
\begin{itemize}
  \item uniform: assign each new peer to the least-loaded pool;
  \item separate pools for peers that send large messages and peers that send small
        ones;
  \item separate pools for unreliable or churning peers.
\end{itemize}

\paragraph{Why the data makes it plausible.}
\begin{itemize}
  \item \textbf{No CPU starvation.} A bounded number of threads avoids the collapse
        at 1,024 peers; the single-poller modes served 1,024 peers at about 140\% CPU.
  \item \textbf{Uses the cores.} One I/O thread plateaus at about 122--144\% CPU in
        S5, while the thread modes reach 170--180\%. Roughly one pool per core could
        reach the same.
  \item \textbf{\code{poll} is cheap for small pools.} Its cost grows with the number
        of descriptors (S2), and the poll-based duplex and split modes had the best
        fan-in latency.
  \item \textbf{Fewer descriptors.} One eventfd per pool instead of one per peer gives
        about \emph{peers + pools}.
\end{itemize}

\paragraph{Preconditions and risks.}
\begin{itemize}
  \item \textbf{The sender's policy has to move into the pool loops.} Retries and
        dials (C3) must become timers there. Otherwise the S4 tail and the S6a cliff
        remain. A plain sleep inside a pool would stall every peer in that pool.
  \item \textbf{Place peers once.} A peer should be assigned when it attaches, without
        migration at first. Moving a live connection means handing over its fd and
        its \code{FrameChannel}.
  \item \textbf{Some isolation policies lack evidence.} Isolating large-message peers
        is motivated by S5 but untested. Isolating slow readers showed no clear need,
        since pooled writes did not stall in S3.
\end{itemize}
Evaluating it would mean sweeping the pool size (1, 8, 32, all) with about one pool
per core, over S1, S4 and S5, under one capacity rule for every mode.

\paragraph{Other directions.} A completion-based backend
(\code{io\_uring}~\cite{iouring}) behind the same \code{IoStack}, and QUIC~\cite{quic},
which is duplex and multiplexed by design.
